\chapter{Make It Yours}\label{ch:yours}

These skills are one person's working style, built during a physics thesis. They are
meant to be taken apart. This chapter says how to adapt them without losing what makes
them useful: rules that live in one place and reach every agent you use.

\section{Start small}

You do not need to use every skill on the first day. Install the whole \code{skills/}
folder, because the skills link to each other, but begin by calling only three:
\code{scientific-mode}, \code{grilling} and \code{unslop}. The first gives each task a
playbook, the second settles decisions before work begins, the third keeps the prose
readable. Work plainly for a while, watch where agents fail, and add a skill or a check
when the same failure shows up twice \cite[Make it yours]{pstack:guide}.

\section{Keep your settings in an overlay}

Do not edit a shared skill to record where your files are or which approver signs a
cleanup. Keep those facts in one personal file next to your project's agent instructions,
for example \code{PERSONAL.md}, and let the skills refer to ``the user'' and ``the
project''. Table~\ref{tab:yours-owners} shows where each kind of rule belongs.

\begin{table}[htbp]
  \small
  \begin{tabularx}{\textwidth}{L{0.3\textwidth}Y}
    \toprule
    Kind of rule & Where it lives \\
    \midrule
    How anyone should work & A shared skill; change it upstream so everyone gets it. \\
    How you work here & The personal overlay: paths, script calls, approvers, delivery. \\
    What the analysis agreed & The analysis repository or its notes: definitions,
    conventions, supervisor decisions. \\
    What happened this week & Session records: \code{BRIEF.md}, \code{STATE.md},
    \code{RESULT.md}. \\
    \bottomrule
  \end{tabularx}
  \caption{Where a rule belongs. A client's own memory feature is read by that client only,
    so it should hold transient context, not rules.}
  \label{tab:yours-owners}
\end{table}

\section{Let lessons reach the file that owns them}

After a task that taught something, run \code{reflect}: it turns the session's lessons
into edits of the skill, overlay or analysis file that owns each rule \cite[Extract
  lessons; Apply]{skill:reflect}. When you correct the same mistake a second time, run
\code{correct}. It prefers a structural fix or a check whose error names the fix over
another sentence of instructions, and keeps a table that pairs each rule with what
enforces it \cite[Fix each class at the highest level that works; Keep the rule
  table]{skill:correct}. A rule that nothing enforces is the first to be forgotten.

\section{Choose how much to delegate}

The default here is to decide, record and report (Chapter~\ref{ch:decisions}): agents
proceed on reversible work and on defensible scientific defaults, and every default stays
visible until you review it. Which actions count as reversible depends on your project.
If your branches are shared, or your review inbox is a shared folder, say so in your
overlay and the gates move accordingly \cite[Decide, record, report]{skill:actions}.

\section{Stay current without being dragged along}

You can install the repository as a plugin and follow it, or copy the skills you want and
own them. If you copy, record the commit you copied from. This repository does the same
with the skills it adapted from others: each records its upstream commit, and a weekly
job opens an issue when an upstream changes. Adopting a change is then a reviewed edit,
not an automatic update.
